% ---- shared definitions for the homework sheets (input after ../lecture_notes/preamble) ----
\newcounter{hwq}
\newenvironment{question}[2]{%
  \refstepcounter{hwq}\needspace{5\baselineskip}%
  \begin{exbox}\textbf{\textcolor{qblue}{שאלה \thehwq\ --- #1}}\hfill\textbf{(#2 נק')}\par\smallskip}%
  {\end{exbox}}
\newenvironment{parts}{\begin{enumerate}[label=(\alph*),itemsep=3pt]}{\end{enumerate}}
\newcommand{\hwtitle}[3]{%
  \noindent\colorbox{qblue}{\parbox{\dimexpr\textwidth-2\fboxsep}{%
   \color{white}\bfseries\Large תרגיל בית #1: #2\\[3pt]\normalsize\mdseries #3}}\par\medskip}
\newcommand{\hwrules}[1]{%
\begin{goalbox}\small\textbf{הנחיות.}
\begin{itemize}
\item ארבע שאלות, 25 נקודות כל אחת. השאלות בלתי תלויות זו בזו.
\item אלא אם נאמר אחרת: $\LCDM$ שטוח עם $h=0.674$, $\Omega_m=0.315$, $\Omega_\Lambda=1-\Omega_m$ (קרינה מוזנחת), ו-$c/H_0=4448\,\Mpc$, $1/H_0=14.51\,\Gyr$.
\item רוב הסעיפים אנליטיים: גזירה, ואז הצבת מספרים במחשבון. רק סעיפים המסומנים \textbf{[מחשב]} דורשים קוד; בהם צרפו את הקוד (\en{Python} מומלץ), סמנו צירים ויחידות, ושמרו את ה-\en{seed}.
\item כל טענה, נתון או ערך שלא נלמדו בהרצאה --- עם ציטוט (מחבר, שנה, כתב עת).
#1
\end{itemize}\end{goalbox}}
\pagestyle{fancy}\fancyhf{}
\fancyhead[L]{\small\hbox{\textdir TLT\relax\thepage}}
